\subsection{Materiais}

% Plataforma
A plataforma de processamento é responsável pela execução do sistema 
operacional embarcado, pela aquisição dos sinais fisiológicos, pelo 
processamento dos dados coletados, pelo cálculo do escore NEWS2 e pela 
execução dos modelos de aprendizado de máquina empregados na predição de 
deterioração clínica. Com isso, foi escolhida a Raspberry Pi 3B+, uma vez que possui ampla documentação para o Buildroot, comunidade ativa, memória RAM de 1 GB, arquitetura 
quad-core de 64 bits a 1,4 GHz e disponibilidade nativa de interfaces
I\textsuperscript{2}C, SPI, HDMI e DSI.

% Sensoriamento
A partir dessa plataforma, a extração dos parâmetros do NEWS2 é realizada por dois sensores. O sensor PPG MAX30102 reúne LEDs vermelho e infravermelho, que viabilizam a extração da SpO\textsubscript{2}, em um encapsulamento compacto, com ADC de 18 bits e baixo custo. Sua interface I\textsuperscript{2}C é nativamente suportada pelo \textit{kernel} Linux, dispensando \textit{frameworks} de terceiros e viabilizando a implementação de um \textit{driver} dedicado em C++. Já o sensor de temperatura DS18B20 oferece resolução configurável de até 12 bits, correspondente a incrementos de 0,0625°C, e precisão de ±0,5°C na faixa de operação corporal. Sua interface digital 1-Wire, também suportada nativamente pelo \textit{kernel}, dispensa ADC externo e permite a leitura da temperatura diretamente via sistema de arquivos, simplificando a integração ao sistema embarcado.

% O MAX30102 se destaca entre os analisados por reunir LEDs vermelho e infravermelho - que viabilizam a extração da SpO\textsubscript{2} - em um encapsulamento compacto, ADC de 18 bits e baixo custo. Além disso, sua interface I\textsuperscript{2}C é nativamente suportada pelo kernel Linux, dispensando \textit{frameworks} de terceiros e viabilizando a implementação de um \textit{driver} dedicado em C++.

% O DS18B20 se destaca entre os analisados por apresentar resolução configurável de até 12 bits, correspondendo a incrementos de 0,0625°C, e precisão de ±0,5°C na faixa de operação corporal. Sua interface digital 1-Wire, nativamente suportada pelo kernel Linux por meio do driver \texttt{w1-therm}, dispensa ADC externo e permite a leitura da temperatura diretamente via sistema de arquivos, simplificando a integração ao sistema embarcado.

% Perifericos
Para o \textit{feedback} visual, optou-se pelo \textit{display} Waveshare de 7 polegadas, cujo tamanho de tela garante alta legibilidade à beira do leito, com compatibilidade direta via HDMI e suporte a toque capacitivo. Complementando o sistema de alertas, o \textit{buzzer} piezoelétrico emite sinais sonoros quando o escore NEWS2 atinge limiares de risco elevado, enquanto o LED RGB reproduz a codificação de cores do protocolo de forma perceptível à distância. Ambos operam via GPIO, sem circuitos de controle adicionais além de resistores limitadores de corrente.

% Como forma de dar \textit{feedback} imediato foi escolhido o \textit{display} Waveshare de 7 polegadas, uma vez que o seu tamanho de tela garante alta legibilidade e usabilidade à beira do leito, e possui compatibilidade direta via interface HDMI com a placa de processamento escolhida e ao suporte a toque capacitivo, que possibilita uma interação mais responsiva com a interface gráfica.

% O buzzer piezoelétrico e o LED RGB complementam o sistema de alertas de forma redundante e de baixo custo. O buzzer emite sinais sonoros quando o escore NEWS2 atinge limiares de risco elevado, garantindo que a equipe de saúde seja notificada mesmo sem visualizar o \textit{display}. O LED RGB, por sua vez, reproduz a codificação de cores do protocolo NEWS2 de forma sintética e imediatamente perceptível à distância, dispensando a leitura ativa do \textit{display} em situações de triagem rápida. Ambos os componentes operam via GPIO, sem necessidade de circuitos de controle adicionais além de resistores limitadores de corrente.

% Alimentação
Quanto à alimentação, optou-se por uma arquitetura conectada diretamente à rede elétrica, o que elimina limitações de autonomia energética e reduz a complexidade do projeto ao dispensar circuitos de gerenciamento de bateria. Para o protótipo, previu-se uma fonte chaveada AC-DC de 5 V, dimensionada para o consumo agregado do sistema e distribuída por um barramento que alimenta em paralelo a plataforma de processamento, o \textit{display} e os demais componentes. A adoção de uma fonte comercial simplifica a implementação do \textit{hardware} e contribui para a segurança operacional, alinhando-se às boas práticas de segurança elétrica para equipamentos eletromédicos estabelecidas pela norma IEC 60601-1 da \textit{International Electrotechnical Commission} (IEC) \citeonline{iec60601}.

% Diante desses requisitos, optou-se por uma arquitetura alimentada diretamente pela rede elétrica. Essa abordagem elimina limitações relacionadas à autonomia energética e reduz a complexidade do projeto ao dispensar circuitos dedicados ao carregamento e gerenciamento de baterias. Além disso, permite maior previsibilidade operacional e simplifica a manutenção do dispositivo.

% Para o protótipo, prevê-se a utilização de uma fonte chaveada AC-DC de 5 V e 3 A, responsável por alimentar a plataforma de processamento, os sensores, o display e os demais periféricos do sistema. A adoção de uma fonte comercial simplifica a implementação do hardware e contribui para a segurança operacional do equipamento, estando alinhada às boas práticas de segurança elétrica aplicadas a equipamentos eletromédicos, conforme estabelecido pela norma IEC 60601-1 da \textit{International Electrotechnical Commission} (IEC) \citeonline{iec60601}.

% Elementos mecanicos
Para a modelagem tridimensional da \textit{case} do dispositivo, utilizou-se o \textit{software} Autodesk Fusion, que dispõe de ferramenta CAD (do inglês \textit{computer-aided design}) com recursos de modelagem paramétrica e licenciamento gratuito para fins educacionais. A fabricação do invólucro foi realizada por manufatura aditiva pelo processo de Modelagem por Deposição Fundida (FDM, do inglês \textit{Fused Deposition Modeling}), em razão de seu baixo custo e facilidade de prototipagem, características que favorecem o desenvolvimento iterativo.

Quanto aos materiais, o Ácido Polilático (PLA, do inglês \textit{Polylactic Acid}) foi adotado nas iterações iniciais, em razão de sua facilidade de impressão e baixo custo, enquanto a versão final empregou o Polietileno Tereftalato Glicol (PETG, do inglês \textit{Polyethylene Terephthalate Glycol}), justificado por sua maior resistência térmica e tolerância à limpeza com álcool isopropílico, que é procedimento de higienização padrão em ambientes hospitalares.

% A modelagem tridimensional da \textit{case} do dispositivo foi utilizado o software Autodesk Fusion, o qual possui a ferramenta de projeto assistido por computador (CAD, do inglês \textit{computer-aided design}, que oferece recursos de modelagem paramétrica, colaboração em nuvem e licenciamento gratuito para fins educacionais e de desenvolvimento. Para a fabricação do invólucro do protótipo foi adotada a manufatura aditiva pelo processo de Modelagem por Deposição Fundida (FDM, do inglês \textit{Fused Deposition Modeling}) devido ao seu baixo custo de implementação, rapidez na produção de peças e facilidade de prototipagem, características que favorecem o desenvolvimento iterativo do produto.

% Para os materiais termoplásticos, foram utilizados dois materiais para etapas distintas do desenvolvimento. O Ácido Polilático (PLA, do inglês \textit{Polylactic Acid}) foi utilizado nas iterações iniciais do invólucro, em razão de sua facilidade de impressão, boa estabilidade dimensional e baixo custo.

% Para a versão final do protótipo, foi adotado o Polietileno Tereftalato Glicol (PETG, do inglês \textit{Polyethylene Terephthalate Glycol}). Essa é justificada por duas propriedades críticas para o contexto de aplicação do dispositivo: maior resistência térmica em relação ao PLA e tolerância à limpeza com álcool isopropílico, procedimento de higienização padrão em ambientes hospitalares e clínicos.


% Ferramentas de software
No âmbito do \textit{software}, o sistema embarcado, responsável pela comunicação com os sensores, pelo processamento dos sinais e pelo cálculo do NEWS2, foi desenvolvido em C++, escolhido por seu desempenho e controle fino sobre memória, determinantes em ambientes com recursos computacionais limitados \cite{knuttila2012functional}. O Buildroot foi selecionado como \textit{toolchain} de \textit{cross-compilação} para a construção de um sistema operacional personalizável baseado em Linux, em razão do desempenho demonstrado em trabalhos correlatos \cite{ozturk2023}.

Para os algoritmos de aprendizado de máquina, adotou-se a linguagem Python, consolidada nesse domínio por seu ecossistema robusto de bibliotecas especializadas e por favorecer a prototipagem rápida. Por fim, a interface gráfica embarcada foi desenvolvida com o \textit{framework} Qt, mantendo o C++ já adotado no restante do sistema embarcado, combinação madura e amplamente utilizada em sistemas embarcados \cite{qtAbout}.

% O Buildroot foi selecionado como toolchain de cross-compilação para o desenvolvimento de um sistema operacional personalizável baseado em Linux, devido ao desempenho demonstrado em trabalhos correlatos, como o de \citeonline{ozturk2023}. Para viabilizar o projeto de uma Placa de Circuito Impresso (PCB, do inglês \textit{Printed Circuit Board}), foi selecionado o Autodesk Fusion, uma vez que possui a ferramenta de Automação de Projetos Eletrônicos (EDA, do inglês \textit{Electronic Design Automation}), além de possuir uma curva de aprendizado rápida, interface intuitiva, licenciamento gratuito e um desenvolvimento colaborativo em nuvem.

% Para o desenvolvimento dos algoritmos de ML, foi adotada a linguagem de programação Python por ser amplamente consolidada nesse domínio por dispor de um ecossistema robusto de bibliotecas especializadas, que oferecem implementações otimizadas de modelos, pipelines de treinamento e ferramentas de avaliação. Ademais, a sintaxe simples e expressiva do Python favorece a prototipagem rápida e a legibilidade do código, aspectos relevantes em projetos de pesquisa que demandam iterações frequentes sobre os modelos desenvolvidos.

% Por fim, para o desenvolvimento da interface gráfica embarcada executada diretamente no hardware, foi selecionado o \textit{framework} Qt em conjunto com a linguagem C++. O Qt é uma plataforma madura e amplamente utilizada no desenvolvimento de interfaces gráficas para sistemas embarcados, oferecendo suporte a renderização eficiente, portabilidade entre diferentes arquiteturas e integração nativa com o sistema operacional Linux, adotado neste trabalho \cite{qtAbout}. A escolha do C++ como linguagem base justifica-se pelo seu desempenho superior em ambientes com recursos computacionais limitados, onde o controle fino sobre memória e tempo de execução é determinante para a responsividade da interface \cite{knuttila2012functional}.